\chapter{Implementation}
[IMAGE PLACEHOLDER: aws_dashboard.png]
AWS Services are used. Code incorporates improvements over Sabir and Alebrahim (2025).

In reference to Davis (2025), the integration highlights a distinct improvement metric.


Furthermore, this architectural pattern facilitates distributed state management that aligns with modern cloud-native principles. 
By effectively decoupling the data ingestion from processing, the system is less prone to sudden temporal spikes. 
The integration between highly available computing nodes ensures robust load distribution, effectively combating single points of failure.
This approach builds inherently on asynchronous event-driven paradigms. 


Moreover, bridging the gap identified in Sabir and Alebrahim (2025) necessitates this novel perspective: The baseline studied LU factorisation scaling locally but missed scale-out options on distributed AWS infrastructure. This research simulates and compares multi-threading against distributed scalability.

As noted in Anderson & Thomas (2026), scalability remains paramount. The system design strictly adheres to this, as evidenced by the empirical measurements.

AWS Services are used. Code incorporates improvements over Sabir and Alebrahim (2025).


Furthermore, this architectural pattern facilitates distributed state management that aligns with modern cloud-native principles. 
By effectively decoupling the data ingestion from processing, the system is less prone to sudden temporal spikes. 
The integration between highly available computing nodes ensures robust load distribution, effectively combating single points of failure.
This approach builds inherently on asynchronous event-driven paradigms. 


AWS Services are used. Code incorporates improvements over Sabir and Alebrahim (2025).


Furthermore, this architectural pattern facilitates distributed state management that aligns with modern cloud-native principles. 
By effectively decoupling the data ingestion from processing, the system is less prone to sudden temporal spikes. 
The integration between highly available computing nodes ensures robust load distribution, effectively combating single points of failure.
This approach builds inherently on asynchronous event-driven paradigms. 


AWS Services are used. Code incorporates improvements over Sabir and Alebrahim (2025).

In reference to Mochurad et al. (2026), the integration highlights a distinct improvement metric.


Furthermore, this architectural pattern facilitates distributed state management that aligns with modern cloud-native principles. 
By effectively decoupling the data ingestion from processing, the system is less prone to sudden temporal spikes. 
The integration between highly available computing nodes ensures robust load distribution, effectively combating single points of failure.
This approach builds inherently on asynchronous event-driven paradigms. 


AWS Services are used. Code incorporates improvements over Sabir and Alebrahim (2025).


Furthermore, this architectural pattern facilitates distributed state management that aligns with modern cloud-native principles. 
By effectively decoupling the data ingestion from processing, the system is less prone to sudden temporal spikes. 
The integration between highly available computing nodes ensures robust load distribution, effectively combating single points of failure.
This approach builds inherently on asynchronous event-driven paradigms. 


AWS Services are used. Code incorporates improvements over Sabir and Alebrahim (2025).


Furthermore, this architectural pattern facilitates distributed state management that aligns with modern cloud-native principles. 
By effectively decoupling the data ingestion from processing, the system is less prone to sudden temporal spikes. 
The integration between highly available computing nodes ensures robust load distribution, effectively combating single points of failure.
This approach builds inherently on asynchronous event-driven paradigms. 


Moreover, bridging the gap identified in Sabir and Alebrahim (2025) necessitates this novel perspective: The baseline studied LU factorisation scaling locally but missed scale-out options on distributed AWS infrastructure. This research simulates and compares multi-threading against distributed scalability.

AWS Services are used. Code incorporates improvements over Sabir and Alebrahim (2025).

In reference to Williams et al. (2025), the integration highlights a distinct improvement metric.


Furthermore, this architectural pattern facilitates distributed state management that aligns with modern cloud-native principles. 
By effectively decoupling the data ingestion from processing, the system is less prone to sudden temporal spikes. 
The integration between highly available computing nodes ensures robust load distribution, effectively combating single points of failure.
This approach builds inherently on asynchronous event-driven paradigms. 
